\documentclass[]{article}

\usepackage{amsmath}
\usepackage{multirow}
\usepackage{natbib}
\usepackage{graphicx} 


\begin{document}

\subsection{Numerical models and datasets}To investigate the link between Eulerian field statistics and Lagrangian transport, we employed a hierarchy of numerical models generating both spatial velocity fields and tracer particle trajectories.\noindent\textbf{Multifractal Energy Cascade.} For the Eulerian analysis, we used synthetic one-dimensional velocity fields generated by a random multiplicative energy cascade model  \citep{shivamoggi2022generalizedfractaldimensiondissipative}. This model allows for precise control over intermittency through the parameter $\mu$  \citep{arreola2024currentstateatmosphericturbulence}. We analyzed three distinct cases: a non-intermittent, Kolmogorov-like cascade ($\mu=0.0$), a mildly intermittent case ($\mu=0.15$), and a case with realistic intermittency ($\mu=0.28$). These datasets provide the basis for calculating the Eulerian spectral roughness.\noindent\textbf{Kraichnan Model.} To study Lagrangian dynamics, we utilized tracer trajectories from the Kraichnan model, a synthetic three-dimensional velocity field that is Gaussian, homogeneous, isotropic, and delta-correlated in time (white noise). The field is defined by a power-law energy spectrum $E(k) \propto k^{-(1+\zeta_2)}$, which allows the spectral roughness $\xi = 1 + \zeta_2$ to be set as a control parameter  \citep{considera2025transportmultifractalkraichnanflows}. This model isolates the effect of spatial correlations on transport by design, removing any temporal memory effects.\noindent\textbf{Lorenz-96 System.} As a deterministic counterpart, we analyzed the Lorenz-96 model, a system of coupled ordinary differential equations that exhibits spatiotemporal chaos. We used spatial snapshots of the system's state in its chaotic regime ($F=8$) for Eulerian structure function analysis. Additionally, we integrated the trajectories of passive tracer particles advected by the Lorenz-96 velocity field to obtain Lagrangian statistics.\noindent\textbf{One-Dimensional Kolmogorov Turbulence.} To probe the influence of spatial dimensionality, we also analyzed tracer displacement data from a synthetic one-dimensional turbulent flow  \citep{benzi199611dimensionalturbulence,kuksin2021kolmogorovstheoryturbulencerigorous} emulating both pure Kolmogorov ($\mu=0.0$) and intermittent statistics  \citep{drivas2026mathematicaltheoremsturbulence}.\subsection{Eulerian spectral roughness analysis}The primary characteristic of the Eulerian velocity field we sought to quantify is its spectral roughness, $\xi$  \citep{jafari2026statisticalfluxfreezingmagnetic}. This was determined from the scaling of the second-order velocity structure function. For a given one-dimensional velocity field $u(x)$, the $p$-th order structure function is defined as the $p$-th moment of the velocity increments over a spatial lag $r$:\begin{equation}S_p(r) = \langle |u(x+r) - u(x)|^p \rangle\end{equation}where $\langle \cdot \rangle$ denotes a spatial average. In a statistically self-similar field, the structure functions exhibit power-law scaling, $S_p(r) \propto r^{\zeta_p}$, where $\zeta_p$ are the scaling exponents  \citep{eyink2000selfsimilardecaykraichnanmodel}. We extracted these exponents for each dataset by performing a linear regression on the log-log plot of $S_p(r)$ versus $r$ over the inertial range. The spectral roughness exponent $\xi$ is then directly defined from the second-order scaling exponent as:\begin{equation}\xi = 1 + \zeta_2.\end{equation}\subsection{Lagrangian transport analysis}For the Lagrangian datasets, we analyzed the statistical properties of tracer trajectories  \citep{biferale2024turblagrdatabase3dlagrangian}. A key prerequisite for the emergence of fractional diffusion is that the Lagrangian velocity statistics are approximately white-in-time. To verify this, we computed the Lagrangian velocity autocorrelation function for each trajectory ensemble  \citep{dewit2026datadrivenmorizwanzigmodelinglagrangian}:\begin{equation}R_v(\tau) = \langle v(t) \cdot v(t+\tau) \rangle,\end{equation}where $v(t)$ is the velocity of a tracer at time $t$. The integral correlation time, $\tau_c$, was then calculated and compared against the total simulation duration to ensure sufficient temporal scale separation.The transport regime was characterized using two primary metrics. First, we computed the Mean Squared Displacement (MSD), $\langle \Delta x^2(\tau) \rangle$, where $\Delta x(\tau) = x(t+\tau) - x(t)$ is the particle displacement over a time lag $\tau$. The scaling of the MSD, $\langle \Delta x^2(\tau) \rangle \propto \tau^{2H}$, yields the Hurst exponent $H$, which was determined via log-log regression. Second, to directly probe the heavy-tailed nature of the displacement distribution predicted by fractional diffusion theory, we estimated the Lévy stability index $\alpha$. This exponent characterizes the power-law tails of the probability density function (PDF) of displacements, $P(\Delta x, \tau)$. We employed a robust Maximum Likelihood Estimator (MLE) on the empirical characteristic function of the displacements to obtain a reliable estimate of $\alpha$.\subsection{Renormalization group flow of the effective exponent}The core of our methodology was to move beyond a single, global measurement of the transport exponent and instead track its evolution as a function of the observational time scale  \citep{bouchara2015anomaloustransportobservableaverage}. This approach allows us to directly visualize the Renormalization Group (RG) flow of the system from its short-time behavior towards its predicted asymptotic fixed point. We define an effective, time-dependent Lévy exponent, $\alpha(\tau)$, which describes the transport at a specific time lag $\tau$.To compute this, we first calculated the empirical characteristic function of the particle displacements for a range of time lags $\tau$  \citep{kaiser2026limittheoremsempiricaldistribution}:\begin{equation}\phi(k, \tau) = \langle e^{ik\Delta x(\tau)} \rangle.\end{equation}For each $\tau$, we then fitted this empirical function to the theoretical characteristic function of a symmetric $\alpha$-stable distribution, which is given by $\phi_{\alpha, D}(k, \tau) = \exp(-D(\tau)|k|^{\alpha(\tau)})$. This fitting procedure yields the effective exponent $\alpha(\tau)$ and an effective diffusivity $D(\tau)$ as a function of the time lag. By plotting $\alpha(\tau)$ versus $\tau$, we can directly observe whether the system flows from the Gaussian regime ($\alpha \approx 2$) towards the theoretically predicted fixed point, $\alpha = 2/\xi$, and determine the crossover timescale over which this transition occurs.

\end{document}
                